\documentclass[a4paper,11pt]{article}
\usepackage[margin=2.5cm]{geometry}

\title{Shikra 0.1 Specification\\[4pt]
\large A small, fast R7RS-small Scheme compiler}
\author{Johannes M\"uller}
\date{19 November 2026 \\ Draft}

\begin{document}
\maketitle

\noindent\textbf{Status.} This is a design specification. Nothing has been
built or measured yet. All size and speed figures are targets.

\tableofcontents
\newpage

% ---------------------------------------------------------------
\section{Overview}

Shikra is an ahead-of-time compiler that turns a whole R7RS-small program
into one small native executable. Compiled programs contain no interpreter
and no \texttt{eval}, so every call target, type and library is visible at
compile time. That is where the speed and the small binaries come from. A
separate interpreter inside the \texttt{shikra} tool provides a REPL for
development.

\subsection{Goals}

In priority order:

\begin{enumerate}
\item \textbf{Runtime speed.} Loops over fixnums, flonums and vectors
  compile to the C loop you would write by hand.
\item \textbf{Small executables.} A binary contains only the code the
  program can reach. Target: hello world under 30\,KB, statically linked on
  Linux.
\item \textbf{The standard language.} R7RS-small as written, with no dialect
  to learn. Any R7RS book or tutorial applies unchanged.
\item \textbf{A small toolchain in plain C99.} One binary, six commands, and
  no dependency beyond a C compiler. Budget: about 15,000 lines for the
  compiler, 2,000 for the REPL interpreter and 4,000 for the runtime.
\end{enumerate}

\subsection{Non-goals}

\begin{itemize}
\item \texttt{eval} or \texttt{load} in compiled programs.
\item R6RS, R7RS-large or \texttt{syntax-case}.
\item Threads in version 1.
\item A full numeric tower by default.
\end{itemize}

\subsection{Design decisions}

\begin{description}
\item[Implementation language] C99 for the compiler, the interpreter and
  the runtime. It builds anywhere with only a C compiler.
\item[Compilation unit] Whole program, closed world. This gives known
  calls, inlining and tree shaking.
\item[Backend] C99 output, compiled by the system \texttt{cc}. This gives a
  mature optimizer and every target for free.
\item[REPL] A bytecode interpreter inside the tool, never in output
  binaries. Interactive development at no cost to compiled programs.
\item[Values] 64-bit tagged words with 63-bit fixnums. Integer and pointer
  operations stay cheap.
\item[Collector] Mark-region, conservative stack, precise heap. Bump
  allocation with no shadow stack.
\item[Tail calls] Loops for known calls, \texttt{musttail} for the rest.
  Guaranteed, with no trampoline cost.
\item[Continuations] Escape-only \texttt{call/cc} is free; re-entrant
  \texttt{call/cc} copies the stack. Ordinary calls pay nothing.
\item[Numbers] Fixnum and flonum. Overflow raises an error and never wraps.
\item[Safety] Checked by default, with an \texttt{--unsafe} opt-out.
\item[Strings] UTF-8 with an ASCII fast path.
\item[Macros] \texttt{syntax-rules} only, exactly what the report requires.
\end{description}

% ---------------------------------------------------------------
\section{Name}

The shikra (\emph{Accipiter badius}) is a small hawk that hunts in short,
fast sprints, which matches a compiler aiming for small and fast output. The
name works as a command: \texttt{shikra build}.

A quick web search found no programming language called Shikra. Two
near-collisions exist: Shirka, a small hobby language
(\texttt{github.com/hashmal/shirka}), and ``The Shikra'', a USB debugging
tool for JTAG, SPI and UART.

Source files keep the standard extensions: \texttt{.scm} for programs and
\texttt{.sld} for library definitions.

% ---------------------------------------------------------------
\section{Language}

Compiled programs get all of \texttt{(scheme base)} and ten of the fifteen
optional libraries. The three evaluator libraries exist in the REPL only.
Complex numbers and \texttt{(scheme r5rs)} are omitted. R7RS-small requires
only the base library; each other library may be provided whole or left out,
so this subset is conformant.

\subsection{Libraries}

\begin{center}
\begin{tabular}{lll}
\hline
Library & Status & Note \\
\hline
\texttt{(scheme base)} & Full & continuations, records, bytevectors, ports \\
\texttt{(scheme case-lambda)} & Full & resolved at compile time for known calls \\
\texttt{(scheme char)} & Full & ASCII tables unless \texttt{--unicode=full} \\
\texttt{(scheme cxr)} & Full & inlined to chained loads \\
\texttt{(scheme file)} & Full & \\
\texttt{(scheme inexact)} & Full & \\
\texttt{(scheme lazy)} & Full & \texttt{delay-force} runs in constant space \\
\texttt{(scheme process-context)} & Full & \\
\texttt{(scheme read)} & Full & linked only when \texttt{read} is reachable \\
\texttt{(scheme time)} & Full & \\
\texttt{(scheme write)} & Full & datum labels for cycles \\
\texttt{(scheme eval)} & REPL only & \texttt{eval} and \texttt{environment} \\
\texttt{(scheme load)} & REPL only & \texttt{load} \\
\texttt{(scheme repl)} & REPL only & \texttt{interaction-environment} \\
\texttt{(scheme r5rs)} & Omitted & exports \texttt{eval} and environments \\
\texttt{(scheme complex)} & Omitted & no complex numbers \\
\hline
\end{tabular}
\end{center}

\subsection{Implementation restrictions}

The report allows each of these. They are where Shikra trades generality for
speed and size.

\begin{itemize}
\item \textbf{Exact integers are 63-bit fixnums.} Overflow raises an error
  and never wraps. Optional bignums may come later, behind a flag.
\item \textbf{No exact rationals.} \texttt{(/ 1 3)} returns the flonum
  0.333\ldots; \texttt{(/ 6 3)} returns the exact 2.
\item \textbf{Inexact reals are IEEE doubles.}
\item \textbf{Strings are UTF-8.} \texttt{string-ref} is constant time on
  ASCII-only strings and linear otherwise. The report does not require
  constant time.
\item \textbf{Macros are \texttt{syntax-rules} only}, fully hygienic. The
  compiler never runs user code at compile time.
\item \textbf{Literal constants are read-only.} Mutating a quoted list or
  string literal is an error per the report; Shikra places them in read-only
  data.
\end{itemize}

\subsection{Errors}

Where the report says ``it is an error'', Shikra in its default safe mode
detects the situation and raises a catchable error object with the procedure
name. That covers wrong types, out-of-range indexes, wrong argument counts
and fixnum overflow. With \texttt{--unsafe} the remaining checks are removed
and behaviour is undefined, as in C.

\subsection{Feature identifiers}

\texttt{cond-expand} sees \texttt{r7rs}, \texttt{shikra},
\texttt{ieee-float}, the operating system and architecture, and
\texttt{full-unicode} when the Unicode tables are enabled. It never sees
\texttt{ratios} or \texttt{exact-complex}.

% ---------------------------------------------------------------
\section{Compilation model}

Shikra compiles the program and every library it imports as one closed unit,
and nothing can be added at run time. This single rule is what makes a
dynamic language compile like a static one.

\subsection{What the closed world gives}

\begin{itemize}
\item \textbf{Known calls.} A top-level definition that is never
  \texttt{set!} is a constant. Calls to it become direct C calls, argument
  counts are checked at compile time, and small procedures are inlined.
\item \textbf{Primitives stay primitive.} R7RS makes it an error to redefine
  or mutate an imported binding. So \texttt{car} and \texttt{+} imported
  from \texttt{(scheme base)} always mean the built-ins and compile to a few
  instructions.
\item \textbf{Private calling conventions.} If every call site of a
  procedure is known, the compiler can pass unboxed integers and doubles,
  drop unused parameters and skip the closure argument.
\item \textbf{Tree shaking.} Anything unreachable from the program body is
  never emitted, including unused parts of the standard library and the
  runtime.
\item \textbf{Closed record types.} Every constructor, accessor and
  predicate of a record type is known, so field access is one tag check and
  one load.
\end{itemize}

\subsection{What it costs}

\begin{itemize}
\item \textbf{Two engines.} The REPL runs an interpreter with an open top
  level, where redefinition is allowed. Compiled code is closed and
  optimized. Both share the front end and the runtime, and one test suite
  runs under both to keep them in agreement.
\item \textbf{Whole-program rebuilds.} Every build recompiles everything.
  Expanded libraries are cached on disk, and the C compiler dominates build
  time.
\end{itemize}

\subsection{How a program is found}

\texttt{shikra build main.scm} treats \texttt{main.scm} as an R7RS program:
import declarations first, then definitions and expressions.
\texttt{(import (geometry vec))} resolves to \texttt{geometry/vec.sld},
searched in the program's directory, then each \texttt{-I} directory, then
the built-in libraries. There is no package manager and no project file.

% ---------------------------------------------------------------
\section{Compiler pipeline}

The compiler runs eight passes over one intermediate representation and
hands a single C file to the system compiler.

\begin{enumerate}
\item \textbf{Read.} Parse source text into datums that carry file, line and
  column.
\item \textbf{Resolve libraries.} Process \texttt{define-library}, import
  sets, \texttt{include} and \texttt{cond-expand}. Build the library graph
  and order it.
\item \textbf{Expand.} Apply hygienic \texttt{syntax-rules} macros. Lower
  derived forms (\texttt{let}, \texttt{cond}, \texttt{do}, \texttt{guard},
  \texttt{parameterize}, \texttt{define-record-type}, \texttt{case-lambda})
  to a core of \texttt{lambda}, \texttt{if}, \texttt{set!}, \texttt{quote},
  \texttt{letrec*} and calls.
\item \textbf{Normalize.} Rename every variable uniquely. Sort
  \texttt{letrec*} bindings into simple, procedure and complex groups. Box
  variables that are both captured and assigned. Convert to the IR.
\item \textbf{Analyze.} One whole-program flow analysis answers two
  questions: which procedures can reach each call site, and which types can
  reach each variable.
\item \textbf{Optimize.} Inline, fold constants, remove dead code, turn
  tail-recursive procedures into loops, remove closures for procedures
  without free variables, delete proven type checks, unbox fixnums and
  flonums.
\item \textbf{Emit C.} Write static data for literals and symbols, then one
  C function per surviving procedure.
\item \textbf{Compile and link.} Run \texttt{cc} on the generated file
  together with the runtime source, as one translation unit.
\end{enumerate}

Passes 1 to 3 are the front end, 4 to 6 the middle end, and 7 the back end.
The REPL's interpreter branches off after the front end.

\subsection{The IR}

The IR is A-normal form with join points: every intermediate value has a
name, and loops and merges are explicit labelled blocks with parameters. It
is close enough to C that pass 7 is a tree walk, and regular enough that
each optimization is a short rewrite.

\subsection{The type lattice}

The analysis tracks a small set of types: fixnum, flonum, char, boolean,
pair, null, string, symbol, vector, bytevector, procedure, each record type,
and ``any''. It is monovariant, meaning one result per variable and not one
per call path. That keeps compile time roughly linear and the analysis under
about 1,000 lines. Precision comes from inlining small procedures before
analysis.

\subsection{Why C output}

Emitting C gives a production-grade optimizer and every target the C
compiler supports, for a pass of about 1,500 lines. The cost is a dependency
on \texttt{cc}. The IR is designed so a direct arm64 or x86-64 backend can
replace passes 7 and 8 later without touching the rest.

\subsection{Written in C99}

The compiler, the interpreter and the runtime are plain C99 with no
dependencies. Generated code is C99 as well, plus two compiler builtins
hidden behind macros: checked addition and the guaranteed tail call. Both
have portable fallbacks.

\begin{itemize}
\item \textbf{Build.} One command, no configure step, no generated sources,
  no existing Scheme needed to bootstrap:
\begin{verbatim}
cc -std=c99 -O2 src/*.c -o shikra
\end{verbatim}
\item \textbf{Memory.} The compiler is a batch program. It allocates from
  arenas and frees nothing until exit.
\item \textbf{IR nodes.} One tagged union with a kind field. Each pass is a
  \texttt{switch} over the kind.
\item \textbf{Shared code.} The reader and the value representation are
  written once. The compiler, the REPL and any compiled program that calls
  \texttt{read} use the same source.
\item \textbf{Standard library.} Most of it is Scheme source embedded in the
  \texttt{shikra} binary as text and compiled with each program.
\end{itemize}

The cost of C is length: tree rewrites that take five lines with pattern
matching in Scheme take twenty in C. That is why the compiler budget is
15,000 lines.

\subsection{The interpreter}

The REPL shares passes 1 to 3 with the compiler, then compiles the core
forms to a compact bytecode that a small virtual machine executes. It never
touches the optimizer or the C backend, so each entry evaluates instantly.

\begin{itemize}
\item \textbf{Same runtime.} The interpreter uses the same tagged values,
  collector and primitives as compiled programs, so values behave and print
  identically.
\item \textbf{Explicit stack.} The virtual machine keeps its own value
  stack. Tail calls run in constant space and re-entrant continuations copy
  that stack, with no dependence on the C stack.
\item \textbf{Open top level.} Global definitions live in a table and can be
  redefined at any time, as the report describes for a REPL.
\item \textbf{Always checked.} There is no unsafe mode in the interpreter.
\item \textbf{Never shipped.} The interpreter is part of the \texttt{shikra}
  tool and is not linked into compiled programs.
\end{itemize}

% ---------------------------------------------------------------
\section{Value representation}

Every Scheme value is one 64-bit word whose low bits say what it is. Fixnums
get the cheapest tag because integer arithmetic is the most common
operation.

\begin{center}
\begin{tabular}{lll}
\hline
Low bits & Kind & Payload \\
\hline
\texttt{...0} & Fixnum & 63-bit signed integer, stored as \texttt{n << 1} \\
\texttt{001} & Pair & pointer to two words, no header \\
\texttt{011} & Procedure & pointer to code address plus captured values \\
\texttt{101} & Other heap object & pointer to header plus fields \\
\texttt{111} & Immediate & character, \texttt{\#t}, \texttt{\#f},
  \texttt{()}, end-of-file, unspecified \\
\hline
\end{tabular}
\end{center}

Consequences of this layout:

\begin{itemize}
\item \textbf{Fixnum add, subtract and compare work on tagged words
  directly.} \texttt{(+ a b)} is one machine add and one overflow branch.
  Multiply untags one operand.
\item \textbf{\texttt{pair?} and \texttt{procedure?} need no memory access.}
  The answer is in the pointer. \texttt{car} is a tag check and one load.
\item \textbf{Pairs take 16 bytes.} Leaving out the header saves a third of
  the memory in list-heavy programs.
\end{itemize}

\subsection{Heap objects}

Every object except a pair starts with an 8-byte header: 8 bits of type,
8 bits of flags and a 48-bit length. After the header:

\begin{description}
\item[Vector] Elements, one word each.
\item[Bytevector] Raw bytes.
\item[String] Character count, byte count and a pointer to a UTF-8 buffer.
  An ``ASCII only'' flag sits in the header.
\item[Symbol] Pointer to its name. Symbols in source are static and compared
  by address.
\item[Flonum] One IEEE double.
\item[Record] Pointer to its type descriptor, then fields.
\item[Port, promise, parameter, error object] Fixed fields.
\end{description}

A procedure with no free variables is a static constant and is never
allocated.

\subsection{Flonums}

A flonum in a generic position is a 16-byte heap box. The optimizer removes
the box wherever the type analysis proves a variable always holds a flonum:
locals, loop variables, and arguments and results of known procedures are
plain C doubles. Flonum literals are static.

NaN-boxing was the alternative. It makes every float free but shrinks
fixnums and adds a mask to every pointer dereference. Shikra bets on
integers and pointers, and relies on static unboxing for float loops. The
known weak spot is floats stored in generic vectors or lists, where each one
is boxed; \texttt{(shikra f64vector)} covers that case.

% ---------------------------------------------------------------
\section{Memory management}

The default collector is a non-moving mark-region design in the style of
Immix, with a conservatively scanned stack and a precisely scanned heap. It
puts no cost on ordinary code: no shadow stack, no stack maps and no write
barrier.

\begin{itemize}
\item \textbf{Allocation is a pointer bump.} The heap is split into 32\,KB
  blocks, each divided into lines. The allocator bumps a pointer through
  runs of free lines. The fast path is inlined: add, compare, store.
\item \textbf{Collection marks, then reclaims lines.} Marking starts from
  the roots; any line with no marked object becomes free again. Mark bits
  live in a side table, so headerless pairs work.
\item \textbf{The stack is scanned conservatively.} Generated C keeps Scheme
  values in ordinary locals and registers. Any stack word that looks like a
  tagged pointer to the start of a live object keeps that object alive.
\item \textbf{The heap is scanned precisely.} Headers and the pair tag tell
  the collector exactly which fields are values.
\item \textbf{Static data is never scanned or freed.} Literals, symbols and
  closure-free procedures live in the executable's data sections.
\end{itemize}

The collector should fit in about 1,200 lines of C. It is single-threaded in
version 1.

\subsection{No-collector mode}

\texttt{--gc=bump} replaces the collector with a bump allocator that never
frees. It suits short-running command-line tools, removes the collector from
the binary, and makes allocation as cheap as it can be.

\subsection{Possible later additions}

\begin{itemize}
\item Stack allocation of objects that escape analysis proves do not outlive
  their procedure.
\item A generational mode using sticky mark bits. It adds a write barrier,
  so it stays optional.
\end{itemize}

% ---------------------------------------------------------------
\section{Control flow}

Shikra runs Scheme on the native C stack with ordinary calls and returns,
and pays for continuations only where a program really uses them.

\subsection{Tail calls}

Proper tail calls are guaranteed, as the report requires, in two ways:

\begin{itemize}
\item \textbf{Loops.} Self-recursion and mutual recursion among known
  procedures compile to jumps inside one C function. A named \texttt{let}
  becomes a C loop.
\item \textbf{\texttt{musttail}.} Every other tail call uses the compiler's
  guaranteed tail-call attribute, supported by clang and recent GCC. These
  procedures share one C signature, with further arguments in a buffer:
\begin{verbatim}
obj f(obj self, intptr_t argc, obj a0, obj a1, obj a2, obj a3);
\end{verbatim}
\end{itemize}

For a C compiler with no such attribute, the backend falls back to a
trampoline, which is correct but slower.

\subsection{Continuations}

Each \texttt{call/cc} site is classified at compile time by what the program
does with the continuation \texttt{k}.

\begin{description}
\item[Tier 1] The program calls \texttt{k} within the same procedure, such
  as an early exit from a loop. Compiled as a jump or a return. No cost.
\item[Tier 2] The program passes \texttt{k} down to callees that only call
  it before \texttt{call/cc} returns. Compiled as a saved stack position and
  a non-local jump. Costs a few words per capture.
\item[Tier 3] The program stores \texttt{k}, returns it, or re-enters it
  later. Compiled as a copy of the live stack into the heap, copied back on
  re-entry. Cost is proportional to stack depth, per capture and per
  re-entry.
\end{description}

Because the standard library is analysed with the program, common patterns
stay in tier 1 or 2. An example is calling \texttt{k} from a
\texttt{lambda} passed to \texttt{for-each}. Generators and coroutines built
on re-entrant continuations land in tier 3 and will be slower than in a
Scheme designed around them. That is the deliberate trade.

\subsection{Exceptions}

\texttt{raise}, \texttt{raise-continuable}, \texttt{error},
\texttt{with-exception-handler} and \texttt{guard} are implemented directly
on a handler stack, not through general \texttt{call/cc}. \texttt{guard}
evaluates its clauses on top of the raising frame and unwinds only when a
clause matches, so no stack is copied. An uncaught error prints its message
and irritants to standard error and exits with a non-zero status.

\subsection{Dynamic state}

\texttt{dynamic-wind} keeps a list of active before and after thunks; any
non-local jump walks that list. \texttt{parameterize} stores its bindings in
the same list, so parameters and unwinding stay consistent.

\subsection{Multiple values}

\texttt{(values a b)} returning to a known \texttt{let-values} or
\texttt{call-with-values} receiver comes back in registers through a small C
struct. Only values flowing to an unknown receiver are packed into a heap
object.

% ---------------------------------------------------------------
\section{Small-binary strategy}

A Shikra executable contains the procedures the program can reach, the
runtime pieces those procedures need, and nothing else. Size is handled as a
reachability problem at three levels.

\begin{enumerate}
\item \textbf{Scheme level.} The compiler emits only definitions reachable
  from the program body. Most of the standard library is written in Scheme,
  so it is shaken and inlined by the same passes as user code.
\item \textbf{Type level.} Generic procedures pull in code for every type.
  The type analysis narrows them: \texttt{(display "hi")} becomes a string
  write, and \texttt{(display n)} with a fixnum \texttt{n} links only the
  integer printer.
\item \textbf{Linker level.} Program and runtime compile as one translation
  unit with one section per function, and the linker drops unreferenced
  sections.
\end{enumerate}

\subsection{Pay only for what you use}

Each runtime piece is linked only under its condition:

\begin{description}
\item[Collector] The program allocates and \texttt{--gc=bump} is not set.
\item[Symbol table] \texttt{string->symbol} or \texttt{read} is reachable.
\item[Float printing and parsing] A flonum can reach a printer, or
  \texttt{string->number} is reachable.
\item[Stack-copying continuations] Some \texttt{call/cc} site is tier 3.
\item[Datum reader] \texttt{read} is reachable.
\item[File ports] \texttt{(scheme file)} procedures are reachable.
\item[Unicode tables] \texttt{--unicode=full} is set.
\end{description}

\subsection{No C library on Linux}

On Linux the runtime makes system calls directly and provides its own entry
point, so the result is a static executable with no libc. On macOS every
executable must link the system library dynamically, and the Mach-O format
page-aligns each segment. Expect a floor of a few tens of kilobytes there,
whatever the program does.

\subsection{Size targets}

\begin{center}
\begin{tabular}{lll}
\hline
Program & Linux x86-64, static & macOS arm64 \\
\hline
Hello world & under 30\,KB & about 50\,KB (format floor) \\
1,000 to 2,000 lines & 60 to 150\,KB & 100 to 200\,KB \\
\hline
\end{tabular}
\end{center}

\texttt{shikra size} prints bytes per library and per procedure, so growth
can be traced to its cause.

% ---------------------------------------------------------------
\section{Extensions}

Shikra adds five small libraries under the \texttt{(shikra ...)} prefix and
one compiler-known form, \texttt{define-c}. Standard programs never need
them; they exist for the cases where the optimizer needs a hint or the
program needs C.

\begin{description}
\item[\texttt{(shikra fixnum)}] \texttt{fx+}, \texttt{fx-}, \texttt{fx*},
  \texttt{fxquotient}, \texttt{fx<?} and the rest of the SRFI 143 names. For
  integer hot paths that must stay fixnum-only.
\item[\texttt{(shikra flonum)}] \texttt{fl+}, \texttt{fl*}, \texttt{flsqrt}
  and the rest of the SRFI 144 names. For float hot paths that must stay
  unboxed.
\item[\texttt{(shikra f64vector)}] Vectors of raw doubles, with SRFI 4
  names. For numeric arrays, signal and matrix code.
\item[\texttt{(shikra assume)}] The form \texttt{(assume expr)} from
  SRFI 145. Checked in safe mode, trusted with \texttt{--unsafe}.
\item[\texttt{(shikra ffi)}] \texttt{define-c}, for calling C functions.
\end{description}

A C binding is one line. Arguments and results are converted to and from
unboxed C values at the call:

\begin{verbatim}
(import (scheme base) (shikra ffi))

(define-c c-hypot "hypot" (double double) double)
(define-c c-getpid "getpid" () int)
\end{verbatim}

Portable code guards these imports with \texttt{cond-expand} on the
\texttt{shikra} feature.

% ---------------------------------------------------------------
\section{Toolchain}

The whole toolchain is one executable, \texttt{shikra}, with six commands.

\begin{description}
\item[\texttt{shikra build main.scm -o app}] Compiles the program to a
  native executable.
\item[\texttt{shikra run main.scm}] Builds at \texttt{-O0} into a cache
  directory and runs it.
\item[\texttt{shikra repl [main.scm]}] Starts the interactive interpreter;
  with a file, loads it first.
\item[\texttt{shikra check main.scm}] Expands and analyses only; reports
  errors without generating code.
\item[\texttt{shikra emit-c main.scm}] Prints the generated C, for reading
  or for an unusual build setup.
\item[\texttt{shikra size app}] Breaks the binary down by library and
  procedure.
\end{description}

\subsection{Flags}

\begin{center}
\begin{tabular}{lll}
\hline
Flag & Values & Default \\
\hline
\texttt{-O} & \texttt{0} fast build, \texttt{2} speed, \texttt{s} size & \texttt{2} \\
\texttt{--unsafe} & removes remaining run-time checks & off \\
\texttt{--gc} & \texttt{immix}, \texttt{bump} & \texttt{immix} \\
\texttt{--unicode} & \texttt{ascii}, \texttt{full} & \texttt{ascii} \\
\texttt{-I dir} & adds a library search directory & none \\
\texttt{--cc} & the C compiler to invoke & \texttt{cc} \\
\texttt{--target} & a target triple for cross-compiling & the host \\
\texttt{-g} & adds source locations to run-time errors & off \\
\hline
\end{tabular}
\end{center}

\subsection{The REPL}

\texttt{shikra repl} is for exploring and learning; \texttt{shikra build} is
for the result. The REPL evaluates each entry with the bytecode interpreter
and prints the value.

\begin{verbatim}
$ shikra repl
shikra> (define (square x) (* x x))
shikra> (map square '(1 2 3))
(1 4 9)
shikra> ,expand (when (> 1 0) 'yes)
(if (> 1 0) (begin 'yes))
\end{verbatim}

\begin{itemize}
\item \textbf{Multi-line input.} An entry continues until its parentheses
  balance.
\item \textbf{Loading a program.} \texttt{shikra repl main.scm} evaluates
  the file, then prompts with its imports and definitions available.
\item \textbf{Errors return to the prompt.} An uncaught error prints its
  message and irritants; the session and its definitions survive.
\item \textbf{Three meta-commands.} \texttt{,expand expr} shows a macro
  expansion, \texttt{,c name} prints the C the compiler would emit for a
  definition, and \texttt{,quit} exits.
\item \textbf{No line-editing library.} History and editing come from a
  terminal wrapper such as \texttt{rlwrap}, which keeps the tool
  dependency-free.
\end{itemize}

Two things to keep in mind. The interpreter is much slower than compiled
output, so measure with \texttt{shikra run}, never in the REPL. And the REPL
accepts things the compiler rejects, such as redefining a procedure or
calling \texttt{eval}; \texttt{shikra check} is the test for whether a
program will build.

\subsection{Diagnostics}

Because the compiler sees the whole program, many mistakes that other
Schemes report at run time are compile-time errors here:

\begin{itemize}
\item An unbound identifier, or an import that no library exports.
\item The wrong number of arguments to a known procedure.
\item A value of a provably wrong type, such as \texttt{(car 5)} or
  \texttt{(vector-ref "abc" 0)}.
\item A \texttt{set!} on an imported binding.
\end{itemize}

Unused imports and definitions are warnings. Every message uses the
\texttt{file:line:column} form that editors understand. Run-time errors in
safe mode name the procedure, and add the source location when built with
\texttt{-g}.

% ---------------------------------------------------------------
\section{Worked example}

A plain R7RS loop should come out as a plain C loop with one overflow check
and no allocation. This is standard Scheme with nothing Shikra-specific in
it:

\begin{verbatim}
(import (scheme base) (scheme write))

(define (sum-to n)
  (let loop ((i 0) (acc 0))
    (if (> i n)
        acc
        (loop (+ i 1) (+ acc i)))))

(display (sum-to 1000000))
(newline)
\end{verbatim}

The C below is a sketch of the intended output, not output from a working
compiler. \texttt{FIX(n)} is \texttt{n << 1}.

\begin{verbatim}
static obj sum_to(obj n) {      /* n proven fixnum at its only call site */
    intptr_t i = FIX(0), acc = FIX(0);
    for (;;) {
        if (i > (intptr_t)n) return (obj)acc;
        if (__builtin_add_overflow(acc, i, &acc)) fix_overflow();
        i += FIX(1);            /* i <= n, so this cannot overflow */
    }
}

int main(void) {
    write_fixnum(STDOUT, sum_to(FIX(1000000)));
    write_char(STDOUT, '\n');
    return 0;
}
\end{verbatim}

How each pass contributed:

\begin{itemize}
\item \textbf{Closed world:} \texttt{sum-to} is never reassigned and has one
  call site, so it is called directly and \texttt{n} is known to be a
  fixnum.
\item \textbf{Loop conversion:} the named \texttt{let} calls itself only in
  tail position, so it becomes a \texttt{for} loop.
\item \textbf{Type analysis:} \texttt{i} and \texttt{acc} start as fixnums
  and are only ever updated by fixnum addition, so every type check
  disappears.
\item \textbf{Tagged arithmetic:} comparison and addition run on the tagged
  words without shifting.
\item \textbf{Tree shaking:} \texttt{display} of a known fixnum links the
  integer printer alone. No collector is linked because nothing allocates.
\end{itemize}

\subsection{Where the speed runs out}

The optimizer cannot help in these cases, and code that depends on them runs
at ordinary dynamic-language speed:

\begin{itemize}
\item Calls through procedures stored in data structures, where the target
  is unknown.
\item Floats kept in generic lists or vectors.
\item \texttt{apply} on lists built at run time.
\item Tier 3 continuations.
\item Generic \texttt{equal?}, \texttt{write} or arithmetic on values whose
  type the analysis could not narrow.
\end{itemize}

% ---------------------------------------------------------------
\section{Roadmap}

Seven milestones take Shikra from hello world to a complete, portable
toolchain, and each ends with a test that either passes or does not. The
first target is arm64 macOS with Apple clang; static Linux follows at M6.

\begin{description}
\item[M0 Hello] Reader, library resolution, \texttt{syntax-rules} expander,
  naive C backend, bump allocator. \emph{Done when} hello world and
  \texttt{fib} build and run.
\item[M1 Language] All of \texttt{(scheme base)}: tail calls, multiple
  values, \texttt{dynamic-wind}, exceptions, tier 1 and 2 continuations,
  strings, bytevectors, ports. \emph{Done when} the \texttt{(scheme base)}
  part of an R7RS test suite passes, for example Chibi's
  \texttt{r7rs-tests}.
\item[M2 Collector] Mark-region collector, static literals and symbols.
  \emph{Done when} a stress mode that collects at every allocation passes
  the same suite.
\item[M3 REPL] Bytecode compiler and virtual machine, \texttt{shikra repl},
  the three evaluator libraries. \emph{Done when} the suite passes in the
  interpreter with output identical to the compiled run.
\item[M4 Optimizer] Flow analysis, inlining, loop conversion, closure
  removal, check elimination, flonum unboxing, tree shaking,
  \texttt{shikra size}. \emph{Done when} size targets are met and the
  \texttt{r7rs-benchmarks} suite runs, with results recorded against Chez
  and Gambit.
\item[M5 Complete] Tier 3 continuations, \texttt{(scheme read)}, cyclic
  \texttt{write}, float printing, the remaining libraries. \emph{Done when}
  the full R7RS test suite passes both compiled and interpreted.
\item[M6 Portable] Static Linux without libc on x86-64 and arm64; optional
  bignums. \emph{Done when} the suite passes on macOS arm64, Linux x86-64
  and Linux arm64.
\end{description}

Shikra is C99 from the first commit, so building it never needs an existing
Scheme. From M3 on the interpreter doubles as a reference: any program whose
compiled and interpreted outputs differ has found a bug in one of the two.

% ---------------------------------------------------------------
\section{Risks}

Five risks could force a redesign.

\begin{itemize}
\item \textbf{Stack-copying continuations.} Tier 3 relies on copying the C
  stack, which the C standard does not define, although several long-lived
  Scheme and Lisp systems have done it. Hardware shadow stacks (Intel CET,
  Arm GCS) reject it. Fallback: convert only the procedures that can be live
  during a tier 3 capture to continuation-passing style.
\item \textbf{Analysis precision.} The speed goal stands or falls with the
  type analysis. If it is too coarse, the checks stay in and Shikra performs
  like any other Scheme-to-C compiler. M4 should get the largest share of
  the effort.
\item \textbf{Two engines drifting apart.} The interpreter and the compiler
  implement the same language twice. A difference between them is a bug
  users will hit when moving from the REPL to a build. Mitigation: shared
  front end and runtime, and one test suite run under both.
\item \textbf{Conservative stack scanning.} It can keep a dead object alive
  and it rules out a moving collector. Both are rare problems on 64-bit
  targets, but fragmentation under long-running workloads is untested.
\item \textbf{Whole-program build time.} Every build recompiles everything,
  and \texttt{cc -O2} on one large file gets slow as programs grow. The
  \texttt{-O0} mode and the expansion cache are the mitigation; splitting
  the C output is the fallback.
\end{itemize}

\end{document}
